\section{L7 CPU Architectures}

\begin{KR}{Recipe: Classify a CPU Memory Architecture}
    shared I+D bus $\to$ Von Neumann; fully separate $\to$ Harvard; split L1 I/D but unified L2/RAM $\to$ \important{Modified Harvard} (e.g.\ A53); DSP separate program/data bus $\to$ Harvard.
\end{KR}

\begin{KR}{Recipe: Identify the Parallelism Mechanism}
    multiple FUs (VLIW) = data; multiple pipelines (superscalar) = instruction; spare slots filled by 2nd thread = SMT; separate cores = multicore/AMP.
\end{KR}

\begin{KR}{Recipe: Set Up AMP}
    same ISA $\to$ SMP/big.LITTLE one OS; different ISA $\to$ separate OS/bare-metal per core + openAMP (\texttt{remoteproc} boots ELF, \texttt{rpmsg}/\texttt{virtio} comms).
\end{KR}

\begin{example2}{Ex.\ S4--7 GP-CPU architectures}
    Von Neumann vs Harvard vs Modified Harvard? Ways to parallelise a CPU?
    \tcblower
    Shared / separate / split-L1+unified-L2. Parallelise via: multicore, multiple FUs (data), multiple pipelines (instruction), heterogeneous PEs.
\end{example2}

\begin{example2}{Ex.\ S5 A53 architecture}
    Separate L1-I/L1-D, unified L2 which architecture?
    \tcblower
    \important{Modified Harvard}: Harvard inside (split L1), Von Neumann outside (unified L2/memory).
\end{example2}

\begin{example2}{Ex.\ S8--10 DSP / VLIW}
    What memory architecture? How does execution work?
    \tcblower
    \important{Harvard} (separate program/data bus keeps MAC fed). \important{VLIW}: compiler packs one op per FU (.M/.D/.L/.S) into a wide word, run in parallel (+ software pipelining). Big code, binary-incompatible.
\end{example2}

\begin{example2}{Ex.\ S11--14 Pipelines}
    Hazards? Super-pipeline vs superscalar? Issue/execute order?
    \tcblower
    Hazards: structural/data/control. \important{Super-pipeline} = finer stages (issue faster); \important{superscalar} = more pipelines (several/cycle). Order: in-order $\to$ OoO exec $\to$ \important{OoO issue+exec}.
\end{example2}

\begin{example2}{Ex.\ S15--18 Co-processors}
    How does a co-processor work and what are its drawbacks?
    \tcblower
    CPU passes unrecognised instructions to it; CPU \emph{stalls} (8087 + 8086 on \texttt{FMUL}); own ISA/clock/mem (NEON). Drawbacks: expensive decoders, integration cost, compiler burden.
\end{example2}

\begin{example2}{Ex.\ S19 Why SMT needs superscalar}
    Why does SMT require a superscalar core?
    \tcblower
    SMT fills \important{idle issue slots} with a 2nd thread; only superscalar cores \emph{have} multiple slots/pipelines per cycle. $\sim$30\% gain (cache-stall dependent); security issues.
\end{example2}

\begin{example2}{Ex.\ S20--25 AMP / openAMP}
    AMP scenarios? Different ISA under one OS? How do cores communicate?
    \tcblower
    Scenarios: diff/same ISA, diff clocks, big.LITTLE, diff/no OS. Different ISA under one OS = \important{No} $\to$ separate OSes + \important{openAMP}: \texttt{remoteproc} boots the ELF, \texttt{rpmsg}/\texttt{virtio} ring buffers carry messages. Zynq: 4$\times$A53 SMP + 2$\times$R5 lockstep + FPGA.
\end{example2}
